\chapter{Random Graphs and the Search for Structure}

\begin{center}\itshape
Join each pair of $n$ vertices with probability $c/n$.\\
If $c<1$, every component has size $O(\log n)$.
\end{center}

The probabilistic method of Chapter~13 proved that an object exists. This chapter asks what a typical object looks like. ``With high probability'' still means that the probability tends to $1$ as $n\to\infty$.

\section{The Erd\H os--R\'enyi graph}

Let $G(n,p)$ have $n$ labelled vertices and each edge present independently with probability $p$.

\begin{theorem}[subcritical components]\label{thm:subcrit}
If $p=c/n$ with $0<c<1$, then with high probability every component of $G(n,p)$ has at most $A\log n$ vertices, for a constant $A=A(c)$.
\end{theorem}
\begin{proof}
A component on exactly $k$ vertices contains a spanning tree, and it has no edge leaving its vertex set. The Prüfer encoding puts the trees on a fixed labelled set of $k$ vertices in bijection with the words of length $k-2$ on that set: delete the smallest leaf and record its neighbour, and reverse the record by attaching the next unused label as a leaf. There are therefore $k^{k-2}$ such trees. Let $N_k$ be the number of vertex sets of size $k$ that span a connected subgraph and have no edge leaving the set. A union bound over those trees gives
\[
 \mathbb E[N_k]\le\binom{n}{k}k^{k-2}p^{k-1}(1-p)^{k(n-k)}.
\]
With $p=c/n$ and $\binom{n}{k}\le(en/k)^k$, and with $(1-p)^{k(n-k)}\le\exp(-ck(1-k/n))$,
\[
 \mathbb E[N_k]\le\frac{n}{ck^2}\bigl(ce^{1-c}\bigr)^k\exp\bigl(ck^2/n\bigr).
\]
The function $\ln c+1-c$ has derivative $1/c-1$ and a maximum value $0$ at $c=1$, so $\rho=ce^{1-c}<1$ for $0<c<1$. For $k\le n^{1/3}$ the exponential factor is at most $2$ once $n$ is large. Thus $\mathbb E[N_k]\le(2n/(ck^2))\rho^k$. The sum of these bounds from $k=\lceil A\ln n\rceil$ to $n^{1/3}$ tends to $0$ when $A>1/\ln(1/\rho)$. This does \emph{not} yet rule out components larger than $n^{1/3}$. The following separate exploration argument handles every size.

Explore the component of a fixed vertex by breadth-first search. When a vertex is processed, reveal its edges to as-yet undiscovered vertices. That number of new neighbours is $\operatorname{Bin}(m_i,c/n)$ for some $m_i\le n$. Fresh previously unexamined edges make these binomial draws conditionally independent of the earlier exploration. Add independent dummy Bernoulli trials to make an upper bound $Y_i\sim\operatorname{Bin}(n,c/n)$ at every step, independent across $i$. If the component contains at least $k\ge2$ vertices, its queue cannot empty before $k-1$ vertices have been processed; therefore $\sum_{i=1}^{k-1}Y_i\ge k-1$.

For any $\theta>0$, independence and the elementary exponential Markov inequality give
\[
\begin{aligned}
P(|C(v)|\ge k)&\le e^{-\theta(k-1)}\Bigl(1+\frac cn(e^\theta-1)\Bigr)^{n(k-1)}\\
&\le\exp\bigl(-(k-1)[\theta-c(e^\theta-1)]\bigr).
\end{aligned}
\]
Choose the explicit value $\theta=\ln(1/c)>0$. Then
\[
 a=\theta-c(e^\theta-1)=\ln(1/c)-(1-c)=c-1-\ln c=-\ln\rho>0,
\]
because $\rho=ce^{1-c}<1$. (Differentiating in $\theta$ shows this is also the best exponent furnished by this exponential bound.) A union bound over all $n$ vertices yields
\[
P(\exists\text{ component of size}\ge k)\le n e^{-a(k-1)}.
\]
Taking $k=\lceil A\ln n\rceil$ with $A>1/a=1/(-\ln\rho)$ gives
\[
 P(\exists\text{ component of size}\ge \lceil A\ln n\rceil)
 \le e^a n^{1-aA}\longrightarrow0.
\]
There is no restriction $k\le n^{1/3}$. This completes the theorem. The earlier spanning-tree argument separately provides a useful estimate for a fixed component size; its small-$k$ cutoff was not used to infer anything about larger sizes.
\end{proof}


The constant $A$ depends only on $c$, explicitly through $\rho=ce^{1-c}$. The proof uses both the expected number of isolated connected sets and an independent exploration bound that controls \emph{every} component size. It does not identify a giant component for $c>1$. Exercise~\ref{ex:giant} gives a positive lower bound on the probability that one vertex has a linear component, and stops there. The survival probability of Theorem~\ref{thm:extinction} is the candidate for the fraction of vertices in that component; the limit of the fraction is not proved.

\section{Exercises}

\begin{exercise}[$\star\star$]\label{ex:connectivity}
Let $p=(\ln n+\omega(n))/n$ and assume $0\le p\le1$ for all sufficiently large $n$ (otherwise $G(n,p)$ is not defined). If $\omega\to+\infty$, show that $G(n,p)$ has no isolated vertex with high probability. If $\omega\to-\infty$ while $p>0$, show that an isolated vertex exists with high probability. State separately what happens at $p=0$.
\end{exercise}


\begin{exercise}[$\star\star\star$]\label{ex:giant}
Let $p=c/n$ with $c>1$, and fix a vertex $v$. Show that there exist $\delta>0$ and $\alpha>0$, depending only on $c$, such that $P(|C(v)|\ge\delta n)\ge\alpha$ for every large $n$. Do not claim that a linear component exists with probability tending to $1$: the argument below gives a positive lower bound, not a limit $1$.
\end{exercise}


\section{Hints}

\begin{hint}{ex:connectivity}
The expected number of isolated vertices is $n(1-p)^{n-1}$.
\end{hint}

\begin{hint}{ex:giant}
While fewer than $\delta n$ vertices have been seen, each new vertex has a $\mathrm{Bin}(\lfloor n(1-\delta)\rfloor,p)$ number of fresh neighbours, which is a Galton--Watson offspring law of mean greater than $1$.
\end{hint}


\section{Solutions}

\begin{solution}{ex:connectivity}
\writeup Let $I$ be the number of isolated vertices. Then $\mathbb E[I]=n(1-p)^{n-1}\le n\exp(-(n-1)p)$. If $\omega\to+\infty$, this tends to $0$, so Markov's inequality gives $P(I\ge1)\to0$. If $p=0$ for some $n$, every vertex is isolated and the lower-regime conclusion is immediate for those $n$; assume $p>0$ in the calculation that follows. If $\omega\to-\infty$ and $0<p\le1$, the formula for $p$ implies $0<p\le\ln n/n$ for large $n$, so $-\ln(1-p)=p+O(p^2)$ and $n p^2=O((\ln n)^2/n)\to0$, hence
\[
 \ln\mathbb E[I]=\ln n-(n-1)(-\ln(1-p))=-\omega+o(1)\to+\infty.
\] For the second moment, two distinct vertices are both isolated when $2n-3$ specific edges are absent, so $\mathbb E[I(I-1)]=n(n-1)(1-p)^{2n-3}$. Hence
\[
 \frac{\mathbb E[I(I-1)]}{\mathbb E[I]^2}=\Bigl(1-\frac1n\Bigr)(1-p)^{-1}.
\]
The right side equals $1+O(p+1/n)=1+O((\ln n)/n)$. Therefore $\mathrm{Var}(I)/\mathbb E[I]^2\to0$ once $\mathbb E[I]\to\infty$, and Chebyshev's inequality yields $P(I=0)\to0$.
\end{solution}


\begin{solution}{ex:giant}
\idea While the explored set is still a small fraction of the graph, breadth-first search dominates a supercritical Galton--Watson process.
\writeup Fix $\delta>0$ with $\mu=c(1-\delta)>1$, and set $m=\lfloor n(1-\delta)\rfloor$. Breadth-first search from $v$ inspects, at the $i$th discovered vertex, the edges from that vertex into the yet unexplored set. Those edges have not been looked at before, and the edge indicators are independent, so the number of new neighbours is a binomial random variable. As long as fewer than $\delta n$ vertices have been discovered, the unexplored set has size at least $m$, and a $\mathrm{Bin}(m',p)$ count with $m'\ge m$ is a $\mathrm{Bin}(m,p)$ count plus an independent nonnegative remainder. Deleting the remainder couples the search to a Galton--Watson process $Z$ with offspring law $\mathrm{Bin}(m,c/n)$, up to the first time $\delta n$ vertices are seen. If the total progeny $T$ of $Z$ satisfies $T\ge\delta n$, the component of $v$ has at least $\delta n$ vertices.

The probability generating function of the offspring law is $f_n(s)=(1-c/n+(c/n)s)^m$. For $s\in[0,1]$,
\[
 \ln f_n(s)=m\ln\bigl(1-\tfrac{c}{n}(1-s)\bigr)\to -\mu(1-s),
\]
so $f_n(s)\to f(s)=\exp(\mu(s-1))$, uniformly on $[0,1]$ because the convergence of $(1-a/m)^{m}$ type expansions is monotone in the error $O((c/n)^2 m)=O(1/n)$. Let $\eta_n$ be the extinction probability of $Z$, the smallest root of $s=f_n(s)$ in $[0,1]$, and let $\eta$ be the smallest root of $s=f(s)$. Theorem~\ref{thm:extinction} gives $\eta<1$ because the Poisson$(\mu)$ mean is $\mu>1$. Every limit point $L$ of $\eta_n$ satisfies $L=f(L)$. The value $1$ is a root, but it is not the limit: $f_n'(1)=mc/n\to\mu>1$, and $f_n''(s)=m(m-1)(c/n)^2(1-c(1-s)/n)^{m-2}$ is bounded uniformly for $s\in[0,1]$ and large $n$. Taylor's formula at $1$ therefore provides a fixed $\varepsilon>0$ such that $f_n(s)<s$ for $s\in(1-\varepsilon,1)$ for all large $n$. Since $f_n(0)>0$ and $f_n$ is continuous, its smallest fixed point must be at most $1-\varepsilon$. Thus $\eta_n\to\eta$. A process that survives has infinite progeny, so $P(T\ge\delta n)\ge P(T=\infty)=1-\eta_n$. Therefore $P(|C(v)|\ge\delta n)\ge1-\eta_n$, which is at least $(1-\eta)/2$ for every large $n$.
\end{solution}


Volume~5, \emph{Calculation: From Formula to Insight}, takes the generating function of Chapter~3 and asks for a remainder you can write down. One infinite sequence of tosses, one Gaussian limit, one path, and one existence proof by expectation: those are the four shapes this book has. The primes are not yet these coins.

\section*{Sources and further reading}

\begin{enumerate}
\item P.~Erd\H os and A.~R\'enyi, On the evolution of random graphs, \emph{Publ.\ Math.\ Inst.\ Hungar.\ Acad.\ Sci.} \textbf{5} (1960), 17--61. The component transition at $p=1/n$.
\end{enumerate}
